\section{答题时间整体规划与答题技巧}

\subsection{试卷结构速览}

\begin{center}
\begin{tabular}{c|c|c|c}
题型 & 题量 & 分值 & 建议用时 \\ \hline
单选题 & 8 题 & 40 分 & 25 分钟 \\
多选题 & 4 题 & 20 分 & 15 分钟 \\
填空题 & 4 题 & 20 分 & 15 分钟 \\
解答题 17 题 & 1 题 & 10 分 & 10 分钟 \\
解答题 18--20 题 & 3 题 & 36 分 & 30 分钟 \\
解答题 21--22 题 & 2 题 & 24 分 & 25 分钟 \\
\hline
总计 & 22 题 & 150 分 & 120 分钟
\end{tabular}
\end{center}

\subsection{答题顺序策略}

\textbf{原则}：先易后难，先熟后生，遇难则跳。

\textbf{推荐答题顺序}：
\begin{enumerate}
  \item 单选题 1--6 题（基础题，10--12 分钟，确保全对）
  \item 填空题 13--14 题（基础题，5--6 分钟，确保全对）
  \item 解答题 17--18 题（基础+中档，15--18 分钟）
  \item 多选题 9--10 题（中档题，6--8 分钟）
  \item 解答题 19--20 题（中档，18--20 分钟）
  \item 单选题 7--8 题（中档偏难，8--10 分钟）
  \item 多选题 11--12 题（难题，6--8 分钟）
  \item 填空题 15--16 题（难题，8--10 分钟）
  \item 解答题 21 题（解析几何，12--15 分钟）
  \item 解答题 22 题（导数，第一问必拿，第二问尽力）
  \item 剩余时间检查
\end{enumerate}

\textbf{底线}：120 分钟内，基础题（约 90 分）应确保拿到 85 分以上，中档题（约 30 分）力争拿 20 分，难题（约 30 分）争取步骤分。

\subsection{单选题答题技巧}

\begin{enumerate}
  \item \textbf{直接法}：概念清晰、公式可直接套用时使用，注意计算准确
  \item \textbf{特值法}（首选技巧之一）：
    \begin{itemize}
      \item 取特殊值（$0,\pm1,\pm2$）代入验证
      \item 取边界值（定义域端点、特殊情况）
    \end{itemize}
  \item \textbf{排除法}：先排除明显错误的选项，缩小范围
    \begin{itemize}
      \item 奇偶性排除、正负号排除、定义域排除
      \item 量纲/单位分析排除
    \end{itemize}
  \item \textbf{数形结合}：
    \begin{itemize}
      \item 函数零点 $\to$ 画图看交点
      \item 方程根的分布 $\to$ 画函数图象
      \item 线性规划、参数范围 $\to$ 画图更直观
    \end{itemize}
  \item \textbf{选项结构分析}：
    \begin{itemize}
      \item 选项中有"以上都不对"时谨慎
      \item 两个选项互为相反/倒数时可能其一为真
    \end{itemize}
\end{enumerate}

\subsection{多选题答题技巧}

\textbf{计分规则}：全部选对得 5 分，部分选对得 2 分，有错选得 0 分。

\textbf{核心策略——宁少勿滥}：
\begin{itemize}
  \item 只选\textbf{100\% 确定}的选项
  \item 对某个选项把握不大时，\textbf{宁愿不选}（多得 2 分 vs 丢 5 分）
  \item 至少可以确保 2 分，避免因贪心丢光
\end{itemize}

\textbf{常见题型特征}：
\begin{itemize}
  \item 函数性质类（奇偶性、单调性、周期性）：逐项验证，性价比高
  \item 解析几何类：定性判断优先，不必每题都算到底
  \item 立体几何类：多考查线面位置关系判断，可用实物或特殊位置验证
  \item 正误判断类：重点关注边界情况和特例
\end{itemize}

\textbf{排除技巧}：
\begin{itemize}
  \item 举反例是证明选项错误最有效的方法
  \item 对含参数的命题，取特殊参数验证
  \item 命题中有"一定""都""所有"等绝对词的，通常为假
\end{itemize}

\subsection{填空题答题技巧}

\textbf{特点}：无过程分，只看最终结果。

\begin{enumerate}
  \item \textbf{计算务必准确}，草稿纸工整，便于检查
  \item \textbf{注意多解}：
    \begin{itemize}
      \item 三角方程注意周期性，不要漏解
      \item 含绝对值的方程注意分类讨论
      \item 二次项系数含参数时考虑 $a=0$ 的情况
    \end{itemize}
  \item \textbf{定义域检查}：
    \begin{itemize}
      \item 对数真数 $>0$，分母 $\neq 0$，偶次根式 $\ge 0$
      \item 函数定义域优先原则
    \end{itemize}
  \item \textbf{双空题}：每个空独立计分，两空之间不要互相影响
  \item \textbf{答案格式}：
    \begin{itemize}
      \item 结果要化简（$\frac{2}{4}$ 写成 $\frac12$，$\sqrt{12}$ 写成 $2\sqrt{3}$）
      \item 范围用区间表示，注意开闭
      \item 符号不能错（$\pm$、$>$、$\ge$ 看清楚）
    \end{itemize}
\end{enumerate}

\subsection{解答题答题规范}

\textbf{通用要求}：
\begin{itemize}
  \item 先写"解："，分步标号（1）（2）
  \item 关键公式必须写原公式再代入（如正弦定理先写 $\frac{a}{\sin A} = \frac{b}{\sin B}$ 再代入）
  \item 运算过程可以跳步，但关键推导步骤不能省
  \item 最终答案建议单独一行或用方框框出
\end{itemize}

\textbf{各题型规范要点}：

\begin{center}
\begin{tabular}{c|c}
\textbf{题型} & \textbf{踩分关键} \\ \hline
三角 & 正弦定理/余弦定理公式 + 代入 + 结果 \\
数列 & 通项公式推导过程 + 求和公式 + 计算 \\
立体几何 & 建系证明 + 坐标正确 + 法向量 + 公式 \\
概率统计 & 分布列 + 期望公式 + 代入计算 \\
解析几何 & 方程联立 + 韦达定理 + 条件转化 + 化简 \\
函数导数 & 求导正确 + 单调性讨论 + 极值/最值
\end{tabular}
\end{center}

\subsection{时间管理技巧}

\begin{itemize}
  \item \textbf{卡壳 3 分钟原则}：一道题思考超过 3 分钟没有进展，标记后跳过
  \item \textbf{首轮抢分}：先做所有能快速拿分的题
  \item \textbf{不空题}：选择题不会也猜一个，填空题不会也写个最可能的数值
  \item \textbf{大题第一问必拿}：22 题第一问通常是送分题（求导、求函数值等）
  \item \textbf{最后 5 分钟}：检查答题卡填涂是否正确，选择题序号是否对齐
\end{itemize}

\subsection{检查策略}

\begin{enumerate}
  \item \textbf{第 1 遍}（做完选择题后）：快速检查有无看错条件、符号错误
  \item \textbf{第 2 遍}（做完填空题后）：检查计算结果是否合理、定义域、多解
  \item \textbf{第 3 遍}（全部做完后）：重点检查计算量大的题目，重算关键步骤
  \item \textbf{检查顺序}：先检查容易检查的（排除法重做） $\to$ 再检查计算量大的
\end{enumerate}

\textbf{常用检查技巧}：
\begin{itemize}
  \item 代入特值验证（取一个满足条件的简单值代回去看等式是否成立）
  \item 逆运算检查（解 $\to$ 代回原题验证）
  \item 量纲分析（单位是否一致、结果是否在合理范围）
  \item 与相邻题对比（如果两道题数据有联系，检查是否一致）
\end{itemize}
